\documentclass[11pt,a4paper]{article}

% White Paper style for CIEM4210 project reports

% --- Essential Packages ---
\usepackage[utf8]{inputenc}
\usepackage[T1]{fontenc}
\usepackage{inter}
\usepackage{charter}
\usepackage[margin=2.5cm]{geometry}
\usepackage{graphicx}
\usepackage{booktabs}
\usepackage{amsmath,amssymb}
\usepackage{hyperref}
\usepackage[table]{xcolor}
\usepackage{titlesec}
\usepackage{enumitem}
\usepackage{array}
\usepackage{tabularx}

% --- Visual Language ---
\definecolor{TUDelftBlue}{RGB}{0,166,214}
\definecolor{PlaceholderBlue}{RGB}{234,246,251}
\definecolor{PlaceholderBorder}{RGB}{120,175,196}
\hypersetup{colorlinks=true, linkcolor=TUDelftBlue, urlcolor=TUDelftBlue}

\titleformat{\section}{\Large\bfseries\sffamily\color{TUDelftBlue}}{\thesection}{0.75em}{}[{\titlerule[0.5pt]}]
\titleformat{\subsection}{\large\bfseries\sffamily\color{black}}{\thesubsection}{0.75em}{}
\titleformat{\subsubsection}{\normalsize\bfseries\sffamily\color{black}}{\thesubsubsection}{0.75em}{}

\setlength{\parindent}{0pt}
\setlength{\parskip}{0.65em}
\renewcommand{\arraystretch}{1.2}

% --- Page Budget Controls ---
\newcounter{maxreportpages}
\setcounter{maxreportpages}{15}
\newcounter{appendixstartpage}
\newcommand{\setmaxreportpages}[1]{\setcounter{maxreportpages}{#1}}
\newcommand{\pagecheck}[1]{\marginpar{\raggedright\footnotesize\itshape Target: #1 pages}}
\let\oeoriginalappendix\appendix
\renewcommand{\appendix}{%
    \setcounter{appendixstartpage}{\value{page}}%
    \oeoriginalappendix
}
\AtEndDocument{%
    \clearpage
    \ifnum\value{appendixstartpage}>0
        \ifnum\numexpr\value{appendixstartpage}-1\relax>\value{maxreportpages}
            \typeout{==============================================}%
            \typeout{CIEM4210 PAGE LIMIT WARNING: main report exceeds 15 pages.}%
            \typeout{Appendix pages are excluded from this count.}%
            \typeout{Reduce content in the main report and move detail to appendices.}%
            \typeout{==============================================}%
            \PackageWarningNoLine{style}{Main report exceeds \arabic{maxreportpages} pages before appendices. Shorten the report body or move detail to appendices.}%
        \fi
    \else
        \ifnum\value{page}>\value{maxreportpages}
            \typeout{==============================================}%
            \typeout{CIEM4210 PAGE LIMIT WARNING: document exceeds 15 pages.}%
            \typeout{Reduce content in the main report and move detail to appendices.}%
            \typeout{==============================================}%
            \PackageWarningNoLine{style}{Document exceeds \arabic{maxreportpages} pages. Shorten the report or move detail to appendices.}%
        \fi
    \fi
}

% --- Reusable Template Blocks ---
\newcommand{\instructionplaceholder}[2]{%
    \vspace{0.5em}
    \noindent\fcolorbox{PlaceholderBorder}{PlaceholderBlue}{%
        \parbox{\dimexpr\linewidth-2\fboxsep-2\fboxrule\relax}{%
            \textbf{Instructional Placeholder: #1}\\
            #2
        }%
    }%
    \vspace{0.5em}
}


\begin{document}

% --- Title Page ---
\begin{titlepage}
    \centering
    \vspace*{2cm}
    {\huge\bfseries\sffamily CIEM4210: Computational Modelling\par}
    \vspace{0.75cm}
    {\Large\sffamily Project Report Template\par}
    \vspace{1.25cm}
    {\large\itshape Modelling a Floating Wind Turbine (OC3 Phase-IV Spar-Buoy)\par}
    \vspace{2cm}
    \begin{tabular}{ll}
        \textbf{Project title:} & [Insert project title] \\
        \textbf{Group number:} & [Insert group number] \\
        \textbf{Student names:} & [Insert names and student numbers] \\
        \textbf{Submission date:} & 12 June 2026, 17:00 \\
    \end{tabular}
    \vfill
    {\large\sffamily TU Delft -- Hydraulic and Offshore Structures\par}
\end{titlepage}

\pagenumbering{roman}
\setmaxreportpages{15}

\section*{How to Use This Template}
\instructionplaceholder{Template use}{Keep the main report within 15 pages, move long derivations and raw outputs to appendices, and ensure each claim is supported by numerical evidence, verification, or literature comparison.}

\instructionplaceholder{Project scope}{This template is focused on the numerical modelling workflow of CIEM4210. Your report should cover structural modelling of the floating offshore wind turbine, numerical method derivation (FEM), implementation details, and quantitative result analysis.}

\subsection*{Core + Choice Project Architecture}
All groups must complete the \textbf{Core Model} and then include the additional components for the highest grade.
\begin{itemize}
    \item Core Model (mandatory): 2D FOWT structural model (RNA point mass, thin-beam tower, rigid-body spar-buoy, spring/damper mooring from quasi-static analysis), baseline Airy-wave/Morison loading, equivalent constant nacelle thrust, equivalent constant current force, FEM formulation, dynamic response assessment, and reduced-order/modal extension.
    \item Extensions (for maximum grade): coupled axial-bending tower model, geometrically nonlinear mooring strings, flexible-beam floater representation, irregular-wave loading, time-dependent thrust force, and time-dependent current force.
\end{itemize}

\subsection*{Section Boundary Rules}
\instructionplaceholder{What belongs where}{Keep evidence in the correct section: assumptions and scope in Section 1; discrete FEM formulation in Section 2; implementation and algorithmic choices in Section 3; verification (convergence, stability, modal) before interpreted engineering conclusions in Section 4. Results presented without prior verification are not considered final evidence.}

\subsection*{Coding Ownership and AI-Minimisation}
\instructionplaceholder{Required practice}{Document your own implementation decisions. For every major solver block, state: (i) why this method was selected, (ii) one alternative that was rejected, and (iii) one verification check used to validate behaviour. This is required even when AI tools are used for language support or debugging assistance.}

\subsection*{Written Statement: Restricted Use of Generative AI}
The use of Generative AI tools (for example Copilot, ChatGPT, and Gemini) is restricted in this course report for the following activities:
\begin{itemize}
    \item Generating derivations, final model equations, or numerical conclusions without independent understanding and verification.
    \item Producing full report sections that are meant to evidence your own reasoning and technical judgement.
    \item Producing complete code or solver routines for submission without testing, validating, and understanding the implementation.
\end{itemize}
Students remain fully responsible for all submitted content. If AI is used, complete the in-depth acknowledgement section at the end of this report.

\tableofcontents
\newpage

\pagenumbering{arabic}

\section*{Executive Summary}
\instructionplaceholder{Executive summary evidence}{Write this last. In at most one page, summarise model scope, key numerical choices, validation outcomes, dominant dynamic response conclusions, and main limitations.}

% ============================================================
% Section 1: Model Definition
% Project focus: OC3 phase-IV spar-buoy floating wind turbine
% ============================================================

\section{Model Definition}
\pagecheck{3--4}

\subsection{Assumptions and scope boundaries}
\instructionplaceholder{Required evidence}{List physical assumptions and scope limits before presenting equations or implementation choices. State what is included, what is excluded, expected impact of each simplification, and the operational range for which your model is intended.}

\subsection{Geometry and system definition}
\instructionplaceholder{Required evidence}{Define the 2D model geometry and all major components: RNA point mass, tower model, spar-buoy body, and mooring representation. State dimensions, reference coordinates, and parameter sources used for the OC3/NREL 5MW representation. Obtain dimensions from the relevant OC3/NREL 5MW reference documents and justify all modelling choices.}

\subsection{Structural model components}
\instructionplaceholder{Pass-level requirements}{Include at minimum: RNA point mass, thin-beam tower model, rigid-body spar-buoy formulation, and spring/damper mooring system with properties derived from a quasi-static analysis.}

\instructionplaceholder{Extensions for maximum grade}{Describe optional extensions selected by your team, such as coupled axial-bending tower behaviour, geometrically nonlinear mooring strings, or a flexible-beam floater representation. Justify each extension physically and numerically, and explain why it was prioritised over alternatives.}

\subsection{Environmental loading model}
\instructionplaceholder{Wave loading}{Describe the wave loading approach. Baseline: linear wave theory (Airy waves) for regular waves and Morison force (inertia and drag terms, including added mass and added damping). Extension for maximum grade: irregular waves.}

\instructionplaceholder{Wind loading}{Describe the wind loading approach. Baseline: equivalent constant thrust force at the nacelle. Extension for maximum grade: time-dependent thrust force.}

\instructionplaceholder{Current loading}{Describe the current loading approach. Baseline: equivalent constant current force. Extension for maximum grade: time-dependent current force.}

\subsection{Governing equations and model scope}
\instructionplaceholder{Required evidence}{Present the governing equations for each structural component, state variables, degrees of freedom, coupling interfaces, and the assembled system. Keep this section at the continuous-model level only; do not include discretisation details here.}

% ============================================================
% Section 2: Numerical Discretization
% ============================================================

\section{Numerical Discretization}
\pagecheck{3--4}

\instructionplaceholder{Derivation overview}{Derive the discrete formulation using the Finite Element Method for the full structural system, demonstrating method ownership. Keep this section focused on mathematical discretization, not software implementation. Inclusion of terms not discussed in class will be valued positively if their inclusion is properly justified.}

\subsection{Finite Element formulation for the structural system}
\instructionplaceholder{Required evidence}{Describe the discretisation of each structural domain (tower beam, floater rigid body, RNA point mass). Detail interpolation functions, weak form derivation, element stiffness/mass/damping matrices, assembly procedure, and boundary conditions. Justify element types and approximation spaces for the expected dynamic range.}

\subsection{Finite Element space and shape functions}
\instructionplaceholder{Required evidence}{Define the FE function spaces used for each field variable. Present the shape functions explicitly, specify their continuity requirements, and link these choices to the governing equations derived in Section 1.}

\subsection{Weak form and loading terms}
\instructionplaceholder{Required evidence}{Derive the weak form of the governing equations for the assembled system. Show explicitly how wave loading (Morison force), wind thrust, and current force enter the discrete loading vector. If time-dependent loading is implemented, show how it is incorporated.}

\subsection{Boundary conditions}
\instructionplaceholder{Required evidence}{State all essential (Dirichlet) and natural (Neumann) boundary conditions. Explain how mooring spring/damper terms and sea-bed fixity are applied at the element level or as additional degrees of freedom.}

\subsection{Discretization justification and alternatives}
\instructionplaceholder{Required evidence}{For each major discretization choice, state why it was selected, one credible alternative, and the expected trade-off in accuracy, robustness, or computational cost.}

% ============================================================
% Section 3: Numerical Schemes
% ============================================================

\section{Numerical Schemes and Implementation}
\pagecheck{2--3}

\instructionplaceholder{Implementation logic}{Document how the discrete FEM equations are implemented in Python or MATLAB. The reader should be able to understand solver workflow, data flow, and time integration choices from this section alone. You may use the lecture notebooks as a starting point; more efficient implementations will be valued positively.}

\subsection{Time integration and solver settings}
\instructionplaceholder{Required evidence}{State time-stepping methods, stability constraints (e.g.\ Courant number, critical time step), solver tolerances, and convergence criteria. Explain why the selected scheme is appropriate for the structural dynamics of the FOWT system, taking into account the relevant natural frequencies identified in Section 4.}

\subsection{Algorithm design and code structure}
\instructionplaceholder{Required evidence}{Describe the computational pipeline, key functions/modules, and reproducibility setup. Include the sequence of assembly, load application, time stepping, and post-processing. Highlight efficiency choices and any extensions beyond the lecture notebooks.}

\subsection{Implementation ownership checklist}
\instructionplaceholder{Required evidence}{For each major solver component, report: selected approach, one alternative considered, and why the final choice was implemented. Detailed convergence, stability, and modal verification evidence belongs in Section 4.}

% ============================================================
% Section 4: Numerical Verification and Results Interpretation
% ============================================================

\section{Numerical Verification and Results Interpretation}
\pagecheck{4--5}

\instructionplaceholder{Result strategy}{Present evidence in two stages: (i) numerical verification (convergence, stability, and modal checks), then (ii) interpreted engineering results under relevant operating conditions. Do not present final engineering conclusions before verification evidence is shown.}

\subsection{Numerical verification}
\instructionplaceholder{Required evidence}{Summarise verification outcomes for the implemented FEM model: consistency checks, benchmark comparisons against literature data, and error trends that justify the credibility of the final simulations. Validate structural deformations and displacements against published OC3/NREL 5MW reference results.}

\subsection{Convergence and time-step assessment}
\instructionplaceholder{Required evidence}{Provide mesh (spatial) and time-step sensitivity analysis. Quantify numerical uncertainty and justify the selected discretisation levels used for all reported final results.}

\subsection{Modal analysis}
\instructionplaceholder{Required evidence}{Identify natural frequencies and mode shapes of the assembled structural model. Relate dominant modes to expected forced-response characteristics under wave, wind, and current loading. Compare computed natural frequencies with literature values where available.}

\subsection{Dynamic response of the structure}
\instructionplaceholder{Required evidence}{Present the dynamic response at key locations (nacelle, centre of gravity, fairlead, etc.) under the operating conditions specified in the project description. Include time histories, frequency spectra, and maximum response estimates as appropriate. Discuss the physical interpretation of the results.}

\subsection{Reduced-order model response}
\instructionplaceholder{Required evidence}{Apply the reduced-order model using a selection of relevant mode shapes. Compare the reduced-order response against the full-model solution and discuss the accuracy and efficiency of the reduction.}

\subsection{Limitations and improvements}
\instructionplaceholder{Required evidence}{State model limitations and propose concrete improvements to geometry, loading fidelity, boundary conditions, or numerical schemes for future work.}


\section*{References}
\instructionplaceholder{Reference practice}{Include only sources that support your model formulation, numerical choices, validation strategy, and interpretation of results. Prioritise course material, standards, and peer-reviewed literature.}

\newpage
\appendix

\section{Appendix A: Additional Derivations and Model Data}
\instructionplaceholder{Appendix use}{Place expanded derivations, parameter tables, and supplementary equations here when they are too long for the main report body.}

\section{Appendix B: Reproducibility Material}
\instructionplaceholder{Appendix use}{Provide enough details to reproduce your results: solver settings, mesh/time-step tables, script structure, and links to executable files or notebooks.}

\section{Appendix C: Decision Log and Verification Checklist}
\instructionplaceholder{Appendix use}{Track key modelling decisions and their verification evidence. Use this log to show technical ownership and reduce black-box coding practices.}

\begin{table}[h!]
    \centering
    \caption{Decision log template.}
    \begin{tabularx}{\textwidth}{>{\raggedright\arraybackslash}p{0.20\textwidth} >{\raggedright\arraybackslash}X >{\raggedright\arraybackslash}p{0.22\textwidth} >{\raggedright\arraybackslash}p{0.22\textwidth}}
        \toprule
        Decision area & Selected approach and rationale & Alternative considered and rejected & Verification evidence (test/plot/table) \\
        \midrule
        Structural model & [Fill in] & [Fill in] & [Fill in] \\
        Mooring model & [Fill in] & [Fill in] & [Fill in] \\
        Loading model & [Fill in] & [Fill in] & [Fill in] \\
        Time integration & [Fill in] & [Fill in] & [Fill in] \\
        \bottomrule
    \end{tabularx}
\end{table}

\section*{Acknowledgement of AI Use}
\instructionplaceholder{If AI was used}{If you used any generative AI tools, complete this section. State: (i) which tool was used, (ii) for which specific task, (iii) how the output was verified or modified, and (iv) what you learned from the process that you would not have learned otherwise. This section is mandatory if AI was involved and will not affect your grade negatively if completed honestly.}

\end{document}
